
\section{第~10~章的练习}
\label{sec:ch14-exercises-ch10}

\vkexsync{46}%

\begin{exercise}\label{ex:ch14-046-fredholm-index}
考虑算子
\[
\begin{array}{c}
L_0\colon W^{1,2}(\RR;\RR^{N})\longrightarrow L^{2}(\RR;\RR^{N})\\[3pt]
Y\longmapsto\dfrac{dY}{ds}+A_0(s)Y\text{，}
\end{array}
\]
其中 $A_0(s)$（对每个 $s$）是一个对角矩阵\footnote{译注：原书此处算子定义印作 $Y\mapsto\dfrac{dY}{ds}+A_0(s)$；$Y$ 是 $\RR^{N}$ 中的向量而 $A_0(s)$ 是矩阵，二者不能相加，按线性齐次算子之文义应为 $Y\mapsto\dfrac{dY}{ds}+A_0(s)Y$；补上 $Y$ 后，(1) 中由两端对角 $\pm1$ 的计数即得 $\operatorname{ind}L_0=m_+-m_-$（即 $n_--n_+$，两端\trMorse{}\trIndex{}之差），恰与\cref{thm:lu-index-morse-difference} 相吻合；此处径改。——译者注}，它在 $|s|$ 充分大时为常值，且形如
\[
A=\begin{pmatrix}\mathrm{Id}_{m_+} & 0\\[2pt] 0 & -\mathrm{Id}_{n_+}\end{pmatrix}
\ \text{（当 } s\ge M\text{）}
\qquad\text{以及}\qquad
\begin{pmatrix}\mathrm{Id}_{m_-} & 0\\[2pt] 0 & -\mathrm{Id}_{n_-}\end{pmatrix}
\ \text{（当 } s\le -M\text{）。}
\]
\begin{enumerate}[label=(\arabic*)]
\item 验证：$L_0$ 是一个\trFredholmOperator。以 $m^{\pm}$ 与 $n^{\pm}$ 为变元确定它的\trIndex。
\item 从\cref{sec:index-computation}的方法中获得启发，由此给出\cref{thm:lu-index-morse-difference}的另一个证明。
\end{enumerate}
\end{exercise}
